\documentclass[8pt,landscape]{article}

% PACKAGES
\usepackage[letterpaper, margin=0.1in]{geometry}
\usepackage{amsmath, amssymb, geometry}
\usepackage{multicol}
\usepackage{setspace}
\usepackage{sectsty}
\usepackage{verbatim}
\usepackage{graphicx}
\usepackage{xcolor}
\usepackage{listings}
\usepackage{enumitem}
\usepackage{booktabs} % For better tables if needed

% Define a custom color for code highlighting and environment
\definecolor{myred}{RGB}{204, 0, 0} % A strong red
\definecolor{mygray}{RGB}{240, 240, 240} % Light gray for background
\setlist[itemize]{
    noitemsep,  % Removes vertical space between list items
    topsep=0pt,  % Removes space before the list
    partopsep=0pt, % Removes extra space when a list begins mid-paragraph
    parsep=0pt, % Removes space between paragraphs within an item
    leftmargin=*, % Sets the left margin to the minimum necessary
    labelsep=0.5em, % You can slightly adjust the space between the bullet and the text
}
% LISTINGS SETUP for R and Python
\lstset{
  basicstyle=\ttfamily\scriptsize\color{black},
  backgroundcolor=\color{mygray},
  frame=single,
  frameround=tttt,
  framesep=0pt,        % no padding between code and frame
  rulecolor=\color{black!20},
  rulesep=0pt,         % no space between rules
  breaklines=true,
  columns=fullflexible,
  showstringspaces=false,
  commentstyle=\color{gray},
  keywordstyle=\color{blue!80!black},
  stringstyle=\color{myred},
  breakatwhitespace=true,
  xleftmargin=0pt,     % no left margin
  xrightmargin=0pt,    % no right margin
  aboveskip=0pt,       % no space above listing
  belowskip=0pt,       % no space below listing
  abovecaptionskip=0pt,
  belowcaptionskip=0pt,
  framexleftmargin=0pt,
  framexrightmargin=0pt,
  framextopmargin=0pt,
  framexbottommargin=0pt
}

% FORMATTING
\setstretch{0.9} % Reduce line spacing
\setlength{\parindent}{0pt}
\setlength{\parskip}{0pt}

% Reduce section spacing and change color
\makeatletter
\renewcommand{\section}{\@startsection{section}{2}{0pt}%
    {0.1ex}% space before section
    {0.1ex}% space after section
    {\fontsize{8}{9}\bfseries\color{red}}} % section font: 8pt
\makeatother

% Reduce subsection spacing and make subsections blue
\makeatletter
\renewcommand{\subsection}{\@startsection{subsection}{2}{0pt}%
    {0.1ex}% space before subsection
    {0.1ex}% space after subsection
    {\fontsize{8}{9}\bfseries\color{blue}}} % Subsection font: 8pt
\makeatother

% Custom command for red-highlighted code/commands (for in-line use)
\newcommand{\code}[1]{\textcolor{myred}{\texttt{#1}}}

% Custom small text wrapper - ensuring 8pt for body text
\newcommand{\smalltext}[1]{
  {\fontsize{8}{9}\selectfont\sloppy #1\par}
}


% DOCUMENT START
\begin{document}
\fontsize{8}{9}\selectfont 
\pagestyle{empty}
\begin{multicols}{3}

\section{Bayesian Normal Linear Regression}
\subsection{Foundations and Likelihood}
\smalltext{
The goal is to assess associations between a response $Y$ and regressors $X$ or predict new observations $Y'$. In the Bayesian framework, we assume response normality where $Y_i \sim N(\mu_i, \sigma^2)$. $\mu_i$ represents the \textbf{systematic component} of the regression model, and $\sigma^2$ is the common variance shared by all responses.
}
\begin{lstlisting}[language=R]
# OLS vs Bayesian Comparison
# OLS can produce negative intercepts that lack physical sense. 
# Bayesian modeling can fix this via restricted prior support.
\end{lstlisting}

\subsection{Stan Program Blocks and Prior Specifications}
\smalltext{
Stan utilizes specific blocks: \code{data}, \code{parameters}, and \code{model}. The \code{generated quantities} block is executed \textbf{after} obtaining posterior samples and is used for predictive inquiries.

\textbf{Posterior variability in parameters:} is the variabilit from approximate joint posterior distribution. These predictions will be computed using the fixed values \code{pred\_x\_1} and \code{pred\_x\_2}.

\textbf{Sampling variability:} Normal random number generator normal\_rng() is added for random noise in MCMC
}
\smalltext{
\begin{itemize}
    \item \textbf{vector:} Linear-algebra type with operator support; preferred for dot products and matrix multiplication like $X * \beta$.
    \item \textbf{array:} Container type best for indexing, nested structures, and element-by-element loops.
    \item \textbf{int:} Integer data \textbf{must} be declared as arrays.
\end{itemize}
}

\begin{lstlisting}[language=R]
model {#Priors allow for clever modeling by choosing appropriate support for parameters.
  beta_0 ~ gamma(7.5, 1);#for Intercept A Gamma prior is used if the response is naturally non-negative.
  beta_1 ~ normal(0, 1000);#for coeff Typically assigned a Normal prior and A large variance reflects prior uncertainty regarding associations.
  sigma ~ inv_gamma(0.001, 0.001);#The Inverse Gamma distribution is a popular non-negative choice.
  y ~ normal(beta_0 + beta_1 * x_1, sigma);}
\end{lstlisting}

\begin{lstlisting}[language=R]
library(bayesrules)
summarize_beta_binomial(a = 4, b = 12, y = 20, n = 200)
\end{lstlisting}

\section{Bayesian Hypothesis Testing}
\subsection{One-Sided Tests and Odds}
\smalltext{
Unlike frequentist testing, each hypothesis ($H_0$ and $H_a$) is assigned a \textbf{probability} derived from the posterior distribution. 
\begin{itemize}
    \item \textbf{Posterior Odds:} The ratio $P(H_a | Y) / P(H_0 | Y)$.
    \item \textbf{Prior Odds:} The ratio $P(H_a) / P(H_0)$.
\end{itemize}
}
\subsection{The Bayes Factor}
\smalltext{
The \textbf{Bayes Factor} is the substitute for the p-value. It compares posterior odds to prior odds to see how inference changed given evidence.
\textbf{$BF = 1$:} Plausibility of $H_a$ remained the same.
\textbf{$BF > 1$:} Plausibility of $H_a$ increased.
\textbf{$BF < 1$:} Plausibility of $H_a$ decreased.
}
\begin{lstlisting}[language=R]
posterior_prob_H_0 <- pbeta(0.15, 24, 192)
posterior_prob_H_a <- pbeta(0.15, 24, 192, lower.tail = FALSE)
posterior_odds <- posterior_prob_H_a / posterior_prob_H_0
prior_prob_H_0 <- pbeta(0.15, 4, 12)
prior_prob_H_a <- pbeta(0.15, 4, 12, lower.tail = FALSE)
prior_odds <- prior_prob_H_a / prior_prob_H_0 
#1/prior_odds times as plausible as the H_a that this mean exceeds
bayes_factor <- posterior_odds / prior_odds
round(qbeta(c(0.025, 0.975), 24, 192), 2) #2 sided test If the 95% credible interval falls within this range, there is reasonable evidence for H_0.
\end{lstlisting}

\subsection{Ethics: B-Hacking}
\smalltext{
\textbf{B-hacking} involves changing priors after seeing data to make results look decisive. To avoid this, commit to priors beforehand and perform \textbf{prior sensitivity checks} by refitting the model with slightly stronger/weaker priors to ensure stable conclusions.
}

\section{Bayesian Binary Logistic Regression}
\smalltext{
Bayesian regression can be extended to \textbf{GLMs}, specifically for binary responses. The key modeling term is the \textbf{probability of success} ($\pi_i$), assumed to follow a $Bernoulli(\pi_i)$ distribution.
}
\subsection{The Link Function}
\smalltext{
To relate the binary response to regressors, we use the \textbf{link function}, which is the \textbf{log of the odds}. It is defined as: $log(\frac{\pi_i}{1 - \pi_i}) = \beta_0 + \beta_1 X_{i,1} + \dots + \beta_k X_{i,k}$. The systematic component allows parameters to be unrestricted on the real scale while the probability $\pi_i$ remains in $[0-1]$ via the \textbf{sigmoid function}.
}

\subsection{Prior Predictive Distribution}
\smalltext{
A \textbf{prior-only model} uses Monte Carlo simulation (not MCMC) to draw from prior distributions without data. This allows us to visualize \textbf{prior sigmoid functions}, which represent what we believe is plausible before seeing evidence. Prior predictive distributions for probabilities often appear non-sensical or extreme before being updated by the likelihood.
}

\subsection{Stan Program Blocks and Prior Specifications}
\smalltext{
\code{data}: Includes the training size $n$, the vector of predictors, and the array of binary responses.
\code{parameters}: Declares the regression coefficients $\beta_0$ and $\beta_1$.
\code{model}: Specifies the priors and uses \code{bernoulli\_logit()} to indicate the logit link function for the likelihood.
The \code{generated quantities} block creates derived values after sampling. \code{inv\_logit()} is used to transform log-odds back into the probability scale ($\hat{\pi}$). To simulate future binary outcomes (0/1), a \textbf{Bernoulli random number generator} (\code{bernoulli\_rng}) is applied to the predicted probability.
}

\begin{lstlisting}[language=R]
model { and beta_1 (slope) typically receive 
  beta_0 ~ normal(0, 100); #Parameters beta_0 (intercept) and beta_1 (slope) Normal priors
  beta_1 ~ normal(0, 100); #mean of zero ($\mu=0$) is often assumed to reflect prior neutrality regarding associations
  y ~ bernoulli_logit(beta_0 + beta_1 * x); #large variance makes the prior weakly informative
}#bernoulli_logit handles log-odds
generated quantities {  real pi_pred;
  int y_pred;  pi_pred = inv_logit(beta_0 + beta_1 * income_pred);
  y_pred  = bernoulli_rng(pi_pred);}
\end{lstlisting}

\subsection{Inference and Comparison}
\textbf{Credible Intervals vs. P-values|}
\smalltext{
In a Bayesian framework, we do not use \textbf{p-values}. Instead, we use \textbf{posterior credible intervals} to test if a coefficient is different from zero. If the 95\% credible interval does not contain zero, we conclude the regressor is associated with the log-odds of the response.
}

\subsection{Bayesian vs. Frequentist (OLS,GLM)}
\smalltext{ 
\textbf{OLS}|A Non-negative Bayesian intercept is more precise than the negative OLS estimation.
Point estimates for differ greatly, but the Bayesian one is more precise.
Both approaches show considerable variability.

\textbf{GLM}|When using weakly informative priors, Bayesian posterior means are often \textbf{practically equal} to frequentist Maximum Likelihood Estimates (MLE) from \code{glm()}. For example, an estimate for $\beta_1$ of 0.009 implies that for each unit increase in the regressor, the subject is $exp(0.009) \approx 1.01$ times more likely to be a "success".
}

\subsection{Posterior Predictive Inference}
\smalltext{Eg.,
our posterior mean indicates that for each one thousand USD increase in income, a subject is 1.01 times more likely to believe in climate change than not to.
}

\section{Bayesian Hierarchical Modeling}
\smalltext{
Standard Bayesian models involve a likelihood conditioned on parameters with fixed hyperparameters. \textbf{Hierarchical models} make these hyperparameters random variables with their own priors. This allows the model to account for two sources of variability: \textbf{Within-group variability} (differences among observations in a group) and \textbf{Between-group variability} (differences from one group to another). It is the best of complete pooled and non-pooled models.
}

\subsection{Modeling Approaches Comparison}
\smalltext{
\begin{itemize}
    \item \textbf{Complete Pooled:} Combines all data into one pool to estimate an overall probability. This ignores individual group differences, potentially leading to high bias.
    \item \textbf{Non-Pooled:} Estimates parameters for each group separately. While flexible, it produces noisy or "wild" estimates for groups with small sample sizes.
    \item \textbf{Hierarchical:} Assumes each group has its own rate, but these rates are drawn from a shared \textbf{population distribution}.
\end{itemize}
}

\subsection{Borrowing Strength}
\smalltext{
  Key property of hierarchical models. Groups with little data "lean" toward the overall pattern of the population, while data-rich groups remain mostly unchanged. This shared information stabilizes estimates and reduces posterior variance compared to non-pooled models.
}

\textbf{Prior Specifications}
\smalltext{
In a hierarchical rocket success model, success probabilities $\pi_i$ follow a $Beta(a, b)$ distribution. The parameters $a$ and $b$ are assigned \textbf{Gamma priors} (Chi-squared also) (e.g., $Gamma(0.001, 0.001)$), which are continuous and non-negative, reflecting the support required for Beta shape parameters.
}

\subsection{Inference and Prediction}
\textbf{Posterior Results:}
\smalltext{
Hierarchical models often yield more \textbf{precise results} with narrower credible intervals than non-pooled models. The model learns the average population probability by analyzing the posterior means of the hyperparameters ($a$ and $b$) through a theoretical Beta distribution.
}

\textbf{Predicting for New Groups:}
\smalltext{
To predict the probability $\pi_{pred}$ for a entirely \textbf{new group} not in the original dataset, the \code{generated quantities} block uses the \code{beta\_rng(a, b)} function. This prediction reflects the population mean and the best Bayesian estimate possible without specific group covariates.
}

\subsection{Relationship to Frequentist Methods}
\smalltext{
Hierarchical Bayesian models are the counterpart to frequentist \textbf{mixed-effects models}. They allow for generalization to new observations and the ability to use information from the entire dataset to infer group-specific probabilities.
}
\section{MCMC Diagnostics: Visual Tools}
\subsection{Trace Plots and Density Overlays}
\smalltext{
  \textbf{Trace plots} illustrate posterior sampling stability; they should show a flat horizontal pattern with no trends or "stuck" chains. \textbf{Empirical density plots} check consistency across multiple chains; overlaid plots should show similar centers, spreads, and shapes for the same parameter. Large discrepancies in density suggest chains are sampling from different posterior regions.
}

\section{MCMC Diagnostics: Quantitative Metrics}
\subsection{Effective Sample Size ($N_{eff}$)}
\smalltext{
$N_{eff}$ measures the number of \textbf{independent samples} required to achieve an accurate approximation. Because MCMC draws are autocorrelated, $N_{eff}$ is often smaller than the raw number of draws $N$. The ratio $N_{eff}/N$ should ideally be close to 1; a ratio $< 1$ suggests a need for more iterations or thinning.
}

\subsection{Autocorrelation}
\smalltext{
  \textbf{Autocorrelation} measures how strongly a draw relates to previous ones. A desirable pattern is for autocorrelation to \textbf{drop quickly} as lag increases (e.g., before lag 5). Slowly decaying autocorrelation indicates inefficient exploration and reduces the effective sample size.
}

\subsection{Gelman-Rubin Diagnostic ($\hat{R}$)}
\smalltext{
$\hat{R}$ evaluates stability by comparing \textbf{within-chain variability} to \textbf{between-chain variability}. Values of $\hat{R} \approx 1$ are desirable, suggesting chains behave similarly. Values noticeably above 1 (e.g., $> 1.1$) indicate instability and that chains have not yet converged to the target posterior.
}

\section{Diagnostic Remediation}
\smalltext{
If diagnostics show poor mixing, drift, or high autocorrelation:
\begin{itemize}
    \item \textbf{Increase Warmup:} If early iterations are influenced by starting values.
    \item \textbf{Run Longer:} Increase iterations to improve $N_{eff}$ and $\hat{R}$ stability.
    \item \textbf{Re-parameterize:} Revisit priors or model specification if problems persist across all diagnostics.
\end{itemize}
}




\end{multicols}
\end{document}
